\subsection{实覆盖空间的 Galois 理论}
\label{sec:02.2.3}

在关于克莱因曲面基本群一般理论的最后这个小节中，我们说明
$\pi^R_1((X,\sigma);x)$ 在什么意义上是 Galois 群。

\begin{theorem}[实覆盖空间的 Galois 理论]\label{thm:2.2.21}
设 $(X,\sigma)$ 是连通克莱因曲面。$(X,\sigma)$ 的实覆盖空间范畴
$\mathrm{Cov}^R_{(X,\sigma)}$ 等价于离散左 $\pi^R_1((X,\sigma);x)$-集范畴。一对显式的
准逆函子由「$x$ 处纤维」函子
\[
\mathrm{Fib}_x:\ \mathrm{Cov}^R_{(X,\sigma)}\longrightarrow
\pi^R_1((X,\sigma);x)\text{-}\mathbf{Sets},
\qquad \bigl(q:(Y,\tau)\to(X,\sigma)\bigr)\longmapsto q^{-1}(\{x\}),
\]
与函子
\[
Q:\ \pi^R_1((X,\sigma);x)\text{-}\mathbf{Sets}\longrightarrow
\mathrm{Cov}^R_{(X,\sigma)},
\qquad F\longmapsto\bigl(\tilde X\times_{\pi_1(X;x)}F,\,
\pi^R_1((X,\sigma);x)/\pi_1(X;x)\bigr)
\]
给出。
\end{theorem}

\begin{proof}
我们在命题 2.2.9 中已看到：实覆盖 $q:(Y,\tau)\to(X,\sigma)$ 在 $x\in X$ 处的纤维上
有 $\pi^R_1((X,\sigma);x)$ 的作用，且实覆盖的同态给出纤维之间关于此作用的等变映射：
这说明 $\mathrm{Fib}_x$ 确实是 $\mathrm{Cov}^R_{(X,\sigma)}$ 到
$\pi^R_1((X,\sigma);x)$-集的一个函子。

反之，设 $F$ 是赋有群 $\pi^R_1((X,\sigma);x)$ 的左作用的离散集。特别地，
$\pi_1(X;x)\subset\pi^R_1((X,\sigma);x)$ 作用在 $\tilde X\times F$ 上，且由定理 2.2.6，
$\tilde X\times_{\pi_1(X;x)}F$ 是 $X$ 的覆盖。为说明
$\tilde X\times_{\pi_1(X;x)}F$ 具有实结构，考虑任一 $h\in\pi^R_1((X,\sigma);x)$，使
其在增广映射下的像为 $(-1)\in\Z/2\Z$（这等价于选取 $\sigma:X\to X$ 的一个提升
$\tilde\sigma:\tilde X\to\tilde X$；由万有覆盖映射 $p:\tilde X\to X$ 的泛性质，满足
$p\circ\tilde\sigma=\sigma\circ p$ 的提升 $\tilde\sigma$ 总存在，但并不唯一，而且一般
地 $\tilde\sigma^2\ne\Id_{\tilde X}$）。于是映射
\[
\Phi:\ \tilde X\times F\longrightarrow\tilde X\times F,
\qquad (y,v)\longmapsto(y\cdot h^{-1},\,h\cdot v)
\]
把 $\pi_1(X;x)$-轨道送到 $\pi_1(X;x)$-轨道：对任意 $f\in\pi_1(X;x)$，
\[
\begin{aligned}
\Phi(y\cdot f^{-1},\,f\cdot v)
&=\bigl((y\cdot f^{-1})\cdot h^{-1},\,h\cdot(f\cdot v)\bigr)\\
&=\bigl((y\cdot h^{-1})\cdot(hfh^{-1})^{-1},\,(hfh^{-1})\cdot(h\cdot v)\bigr)\\
&\sim(y\cdot h^{-1},\,h\cdot v)=\Phi(y,v),
\end{aligned}
\]
这里用到 $hfh^{-1}\in\pi_1(X;x)\subset\pi^R_1((X,\sigma);x)$。因此 $\Phi$ 诱导映射
$\tilde X\times_{\pi_1(X;x)}F\to\tilde X\times_{\pi_1(X;x)}F$，且因 $h^2\in\pi_1(X;x)$
其平方为恒等。换言之，群
$\pi^R_1((X,\sigma);x)/\pi_1(X;x)\simeq\Z/2\Z$ 作用在
$Y:=\tilde X\times_{\pi_1(X;x)}F$ 上。记 $\tau$ 为该 $\Z/2\Z$ 作用诱导的对合，则
$(Y,\tau)$ 是克莱因曲面。由于 $X$ 上的实
结构 $\sigma$ 正是按上述方式取 $F=\{\mathrm{pt}\}$ 得到的，克莱因曲面 $(Y,\tau)$ 确实是
$(X,\sigma)$ 的实覆盖，其纤维为 $F$。此外，$Y$ 上的实结构 $\tau$ 不依赖于 $h$ 的选取，
因为任意两个选取相差 $\pi_1(X;x)$ 中一个元素。这一构造是函子式的，并且是
$\mathrm{Fib}_x$ 的准逆。
\end{proof}

\subsection*{习题}

\begin{exercise}\label{ex:2.2.22}
设 $V$ 是向量空间，$E$ 是以 $V$ 为平移群的仿射空间。证明 $E$ 的仿射变换群是半直积
$V\rtimes\GL(V)$，其中 $\GL(V)$ 是 $V$ 的线性自同构群。回忆该半直积上的复合律为
$(v,A)\cdot(w,B)=(v+Aw,AB)$。
\end{exercise}

\begin{exercise}\label{ex:2.2.23}
设 $X$ 是黎曼面，$\Aut^{\pm}(X)$ 是 $X$ 到自身的全纯与反全纯双射之集。
\begin{enumerate}[label=(\roman*), leftmargin=2.2em, itemsep=0.15em]
\item 证明 $\Aut^{\pm}(X)$ 是群。
\item 证明：若 $\sigma$ 是 $X$ 上的实结构，则 $\Aut^{\pm}(X)$ 同构于
      $\Aut(X)\rtimes\Z/2\Z$，其中 $\Z/2\Z\simeq\{\Id_X,\sigma\}$ 通过
      $\sigma\cdot f=\sigma\circ f\circ\sigma^{-1}$ 作用在 $\Aut(X)$ 上。
\end{enumerate}
\end{exercise}

\begin{exercise}\label{ex:2.2.24}
设 $G,H$ 是两个群，$\phi:H\to\Aut(G)$ 是群同态（即 $H$ 通过群自同构对 $G$ 的左作用）。
\begin{enumerate}[label=(\roman*), leftmargin=2.2em, itemsep=0.15em]
\item 证明 $\phi$ 诱导群同态 $\phi^{\mathrm{op}}:H^{\mathrm{op}}\to
      \Aut(G^{\mathrm{op}})^{\mathrm{op}}$（即 $H^{\mathrm{op}}$ 通过群自同构对
      $G^{\mathrm{op}}$ 的右作用）。
\item 证明映射 $(G\rtimes H)^{\mathrm{op}}\to H^{\mathrm{op}}\ltimes G^{\mathrm{op}}$，
      $(g,h)\mapsto(h,g)$ 是群同构。
\end{enumerate}
\end{exercise}

\begin{exercise}\label{ex:2.2.25}
考虑把 $\tilde X$ 按通常方式构造为从给定 $x\in X$ 出发的道路的同伦类之集，以及 (2.8)
中对任意 $\gamma\in\pi_1(X;x)$、$\eta\in\tilde X$ 定义的作用
$\eta\curvearrowleft\gamma:=\eta*\gamma$。证明：若 $\eta$ 事实上是 $x$ 处环路的同伦
类，则 (2.5) 中定义的 $\gamma$ 对 $\eta$ 的单值作用 $\gamma\curvearrowright\eta$
确实就是 $\gamma*\eta$（即 $\eta$ 与 $\gamma$ 在 $\pi_1(X;x)$ 中的反向乘积）。
\end{exercise}

\begin{exercise}\label{ex:2.2.26}
例如利用 Szamuely 的书 [18] 补出 2.2.1 小节的全部细节。特别地，证明当函子
$\mathrm{Fib}_x:\mathrm{Cov}_X\to\pi_1(X;x)\text{-}\mathbf{Sets}$ 限制在连通覆盖空间上
时，它被 $\tilde X$ 表出。
\end{exercise}

\begin{exercise}\label{ex:2.2.27}
证明：若 $p:(\tilde X,\tilde x)\to(X,x)$ 是万有覆盖，
$f\in\Aut(\tilde X/X)^{\mathrm{op}}$，则对一切 $y\in p^{-1}(\{x\})\subset\tilde X$ 有
$f\circ t_y=t_{f(y)}$，其中 $t_y$ 与 $t_{f(y)}$ 如命题 2.2.5 中所定义。
\end{exercise}

\begin{exercise}\label{ex:2.2.28}
证明复平面 $\C$ 的反全纯对合，被双全纯变换群
$\Aut(\C)\simeq\{z\mapsto az+b\mid a\in\C^*,\ b\in\C\}$ 中某元素共轭于
$\tau_0:z\mapsto\bar z$。
\end{exercise}

\begin{exercise}\label{ex:2.2.29}
证明 Poincaré 上半平面 $\mathbb{H}=\{z\in\C\mid \operatorname{Im}z>0\}$ 的反全纯对合
被 $\mathbb{H}$ 的双全纯变换群 $\Aut(\mathbb{H})\simeq\mathrm{PSL}(2;\R)$ 中某元素共轭于
$\tau_0:z\mapsto-\bar z$。
\end{exercise}

\begin{exercise}[由 Andrés Jaramillo 建议]\label{ex:2.2.30}
设 $(X,\sigma)$ 是 Real 黎曼面，$\tilde X$ 是其万有覆盖黎曼面。
\begin{enumerate}[label=(\roman*), leftmargin=2.2em, itemsep=0.15em]
\item 证明：存在覆盖 $\sigma$ 且满足 $\tilde\sigma^2=\Id_{\tilde X}$ 的映射
      $\tilde\sigma:\tilde X\to\tilde X$，等价于同伦正合列 (2.17) 存在分裂。
\item 设 $\sigma(x)\ne x$。证明同伦正合列存在分裂，当且仅当存在从 $x$ 到 $\sigma(x)$
      的道路 $\eta$，它在 $\pi^R_1((X,\sigma);x)$ 中是二阶元，即
      $(\sigma\circ\eta)*\eta$ 同伦于 $x$ 处的常道路。
\item 证明：若 $\Fix(\sigma)=\varnothing$，则这不可能发生。
\end{enumerate}
\end{exercise}

\begin{exercise}\label{ex:2.2.31}
证明：若 $x$ 与 $x'$ 落在 $X^\sigma$ 的同一连通分支中，则由 $X^\sigma$ 中从 $x$ 到
$x'$ 的道路（的同伦类）定义的增广的同构
\[
\begin{tikzcd}[column sep=large, row sep=large]
\pi_1(X;x)\rtimes_{\tilde\sigma}\Z/2\Z \arrow[r] \arrow[d, "\simeq"'] & \Z/2\Z \arrow[d, "="]\\
\pi_1(X;x)\rtimes_{\tilde\sigma'}\Z/2\Z \arrow[r] & \Z/2\Z
\end{tikzcd}
\]
其中 $\tilde\sigma$ 与 $\tilde\sigma'$ 分别按引理 2.2.11 用 $x$ 与 $x'$ 定义。
\end{exercise}
